\section{Theory}

In the following sections we will examine the emission of different surfaces and of a cavity radiator.
A reasonable theoretical starting point is the black body radiator. Its spectral density of electromagnetic radiation is given by Planck's law:
\begin{equation}
    I(\nu, T) = \frac{2h\nu^3}{c^2} \cdot \frac{1}{e^{\frac{h\nu}{k T}} - 1}
\end{equation}
Here the T is the temerature of the black body, $\nu$ is the frequency of the emitted radiation, $h$ is Planck's constant, $c$ is the speed of light and $k$ is Boltzmann's constant. 
The total power emitted by a black body per unit area can be calculated by integrating over all frequencies and the solid angle:
\begin{equation}
    \frac{P}{A} = \int_0^\infty I(\nu, T) d\nu \int cos(\theta) d\Omega 
\end{equation}
The cosine is there because the black body is also a Lambertian emitter, meaning that the intensity of the emitted radiation is proportional to the cosine of the angle between the normal of the surface and the direction of emission.
Computing the integral yields the Stefan-Boltzmann law:
\begin{equation}
    P = \sigma A T^4
    \label{eq:SBG}
\end{equation}
Where $\sigma$ is the Stefan-Boltzmann constant:
\begin{equation}
    \sigma = \frac{2\pi^5 k^4}{15 h^3 c^2} \approx 5.67 \cdot 10^{-8} \frac{W}{m^2 K^4}
\end{equation}
Since the black body is an idealized object, Eq. \ref{eq:SBG} is not valid for real objects. Instead, we have to introduce the specific emissivity $\epsilon$ of the object, which is smaller than 1 for real objects and equals 1 for a perfect black body.
The modified Stefan-Boltzmann law for real objects is then:
\begin{equation}
    P = \epsilon \sigma A T^4
    \label{eq:SBGreal}
\end{equation}
For the following measurments it is important to note, that the Leslie's cube also absorbs radiation from the environment and re-emits it with respect to the Kirchhoff's law.
Kirchhoff's law states that the emissivity $\epsilon$ of a surface is equal to its absorptivity $\alpha$.
Therefore the radiation generated only by the temperature of the cube is the difference between the emitted radiation and the absorbed radiation.
$T_0$ is the temperature of the environment, which we will assume to be room temperature.
It follows:
\begin{equation}
    P \propto T^4 - T_0^4
\end{equation}